
\documentclass{article}
\usepackage{colortbl}
\usepackage{makecell}
\usepackage{multirow}
\usepackage{supertabular}

\begin{document}

\newcounter{utterance}

\centering \large Interaction Transcript for game `eqbench', experiment `eq\_bench', episode 1 with reflexion{-}Qwen3.5{-}27B.
\vspace{24pt}

{ \footnotesize  \setcounter{utterance}{1}
\setlength{\tabcolsep}{0pt}
\begin{supertabular}{c@{$\;$}|p{.15\linewidth}@{}p{.15\linewidth}p{.15\linewidth}p{.15\linewidth}p{.15\linewidth}p{.15\linewidth}}
    \# & \multicolumn{2}{c}{Player} && \multicolumn{2}{c}{Game Master} \\
    \hline

    \theutterance \stepcounter{utterance}  
    & & & \multicolumn{4}{p{0.6\linewidth}}{
        \cellcolor[rgb]{0.9,0.9,0.9}{
            \makecell[{{p{\linewidth}}}]{
                \texttt{\tiny{[P1$\langle$GM]}}
                \texttt{You have attempted this same task before and did not succeed. These are the plans you wrote after those attempts.} \\
\texttt{Trial \#1: New plan: My previous attempt failed because I misjudged the primary emotional driver of the dialogue. While I correctly identified "Defensive" as high, I overestimated "Incredulous" and underestimated the nuance of the conflict. Claudia's responses ("Not my world," "testament to our parents' love") are direct rebuttals to Robert's attacks, indicating a strong protective stance rather than disbelief. Robert's threats about "survival" and "bulldozing" are aggressive, but Claudia's tone remains firm and argumentative, not fearful (Intimidated) or pitying (Sympathetic). Next time, I will analyze the specific function of each line: if a character is countering an accusation with a moral or personal justification, "Defensive" should be the dominant score (9{-}10). "Incredulous" should be reserved for expressions of shock or disbelief (e.g., "You can't be serious"), which are absent here, so I will lower that score significantly (0{-}2). I will ensure the scores reflect the hierarchy of emotions: Defensive > Intimidated > Incredulous > Sympathetic, based on the text's evidence of a heated ideological clash rather than a reaction to a surprise or a threat of physical harm.} \\
\\ 
\texttt{Trial \#2: New plan: My previous attempt failed because I assigned non{-}zero scores to emotions that were not supported by the text, specifically "Incredulous" and "Intimidated." Claudia's dialogue consists of firm rebuttals and moral standpoints ("Not my world," "testament to our parents' love"), which indicates she is standing her ground rather than being shocked or fearful. Robert's threats about "survival" are met with counter{-}arguments about "trees," showing she is not cowed. Therefore, "Intimidated" and "Incredulous" should be scored as 0. "Sympathetic" is also 0 as she is not expressing pity but rather anger and conviction. The only valid high{-}intensity emotion is "Defensive," as she is actively protecting her values and her parents' legacy against Robert's attacks. Next time, I will strictly assign 0 to any emotion not explicitly evidenced by the character's tone or specific word choices, reserving high scores only for the primary emotional driver, which in this case is purely "Defensive."} \\
\\ 
\texttt{Trial \#3: New plan: My previous attempt failed because I assigned a score of 0 to "Incredulous," overlooking the rhetorical nature of Claudia's responses. While "Defensive" is indeed the dominant emotion (score 10) as she protects her values, her lines like "Practicality, according to you, means bulldozing everything" and "Your world? You mean this...this sanctuary?" are not just defensive; they express disbelief at Robert's distorted view of reality and his reduction of her sanctuary to mere sentimentality. This rhetorical questioning indicates a level of incredulity that warrants a non{-}zero score (e.g., 4{-}5). "Intimidated" and "Sympathetic" remain 0 as she shows no fear or pity. Next time, I will distinguish between pure defense (protecting a stance) and incredulity (expressing disbelief at the opponent's logic), ensuring that rhetorical challenges to the opponent's worldview are scored appropriately rather than dismissed as purely defensive.} \\
\\ 
\\ 
\\ 
\texttt{Your task is to predict the likely emotional responses of a character in this dialogue:} \\
\\ 
\texttt{Robert: Claudia, you've always been the idealist. But let's be practical for once, shall we?} \\
\texttt{Claudia: Practicality, according to you, means bulldozing everything in sight.} \\
\texttt{Robert: It's called progress, Claudia. It's how the world works.} \\
\texttt{Claudia: Not my world, Robert.} \\
\texttt{Robert: Your world? You mean this...this sanctuary of yours?} \\
\texttt{Claudia: It's more than a sanctuary. It's a testament to our parents' love for nature.} \\
\texttt{Robert: Sentimentality won't pay the bills, Claudia.} \\
\texttt{Claudia: And money won't bring back the trees you've felled, Robert.} \\
\texttt{Robert: This isn't about trees, Claudia. It's about survival. My company's survival.} \\
\texttt{{[}End dialogue{]}} \\
\\ 
\texttt{At the end of this dialogue, Claudia would feel...} \\
\texttt{Sympathetic} \\
\texttt{Incredulous} \\
\texttt{Intimidated} \\
\texttt{Defensive} \\
\\ 
\texttt{Give each of these possible emotions a score from 0{-}10 for the relative intensity that they are likely to be feeling each.} \\
\\ 
\texttt{You must output in the following format, including headings (of course, you should give your own scores), with no additional commentary:} \\
\\ 
\texttt{Sympathetic: <score>} \\
\texttt{Incredulous: <score>} \\
\texttt{Intimidated: <score>} \\
\texttt{Defensive: <score>} \\
\\ 
\\ 
\texttt{{[}End of answer{]}} \\
\\ 
\texttt{Remember: zero is a valid score, meaning they are likely not feeling that emotion. You must score at least one emotion > 0.} \\
\\ 
\texttt{Your answer:} \\
            }
        }
    }
     \\ \\

    \theutterance \stepcounter{utterance}  
    & \multicolumn{4}{p{0.6\linewidth}}{
        \cellcolor[rgb]{0.561,0.737,0.561}{
            \makecell[{{p{\linewidth}}}]{
                \texttt{\tiny{[P1$\rangle$GM]}}
                \texttt{Sympathetic: 0} \\
\texttt{Incredulous: 4} \\
\texttt{Intimidated: 0} \\
\texttt{Defensive: 10} \\
            }
        }
    }
    & & \\ \\

    \theutterance \stepcounter{utterance}  
    & & & \multicolumn{2}{p{0.3\linewidth}}{
        \cellcolor[rgb]{0.9,0.9,0.9}{
            \makecell[{{p{\linewidth}}}]{
                \texttt{\tiny{[GM$|$GM]}}
                \texttt{Sympathetic: 0} \\
\texttt{Incredulous: 4} \\
\texttt{Intimidated: 0} \\
\texttt{Defensive: 10} \\
            }
        }
    }
    & & \\ \\

    \theutterance \stepcounter{utterance}  
    & & & \multicolumn{2}{p{0.3\linewidth}}{
        \cellcolor[rgb]{0.9,0.9,0.9}{
            \makecell[{{p{\linewidth}}}]{
                \texttt{\tiny{[GM$|$GM]}}
                \texttt{Sympathetic: 0} \\
\texttt{Incredulous: 8} \\
\texttt{Intimidated: 3} \\
\texttt{Defensive: 7} \\
            }
        }
    }
    & & \\ \\

    \theutterance \stepcounter{utterance}  
    & & & \multicolumn{2}{p{0.3\linewidth}}{
        \cellcolor[rgb]{0.9,0.9,0.9}{
            \makecell[{{p{\linewidth}}}]{
                \texttt{\tiny{[GM$|$GM]}}
                \texttt{game\_result = LOSE} \\
            }
        }
    }
    & & \\ \\

    \theutterance \stepcounter{utterance}  
    & & \multicolumn{4}{p{0.6\linewidth}}{
        \cellcolor[rgb]{0.95,0.95,0.95}{
            \makecell[{{p{\linewidth}}}]{
                \texttt{\tiny{[P1$\leftrightarrow$P]}}
                \texttt{Post{-}episode Reflexion (not sent to GM):} \\
\texttt{New plan: My previous attempts failed because I oscillated between underestimating and overestimating the "Incredulous" score, failing to pinpoint the precise emotional hierarchy. In Trial 1, I gave "Incredulous" a 4, which was likely too high for a character who is firmly arguing rather than expressing shock. In Trial 3, I justified a non{-}zero score for "Incredulous" based on rhetorical questions, but the prompt's outcome suggests the target state requires a stricter interpretation where "Defensive" is the *sole* dominant emotion. Claudia's lines are direct counter{-}attacks ("Not my world," "money won't bring back the trees") rather than expressions of disbelief. She is not surprised by Robert's stance; she is actively resisting it. Therefore, the correct strategy is to assign the maximum score (10) to "Defensive" and strictly 0 to all other emotions, including "Incredulous," as her tone is one of conviction and rebuttal, not shock or disbelief. I will ensure the output reflects this singular emotional focus: Defensive: 10, and all others: 0.} \\
            }
        }
    }
    & \\ \\

\end{supertabular}
}

\end{document}
